% ================================================================================================ capítulo 10
\chapter{Medir los puntos y confiar en los números}\label{cap:estadistica}

\begin{resumen}
\textbf{En este capítulo:} cómo encuentra TurbuLENS los dos puntos de la estrella en cada cuadro, qué estadística usa
para resumirlos, por qué descarta los datos atípicos con la mediana y la MAD, y cómo leer los tres controles de calidad
de \menu{Resultados}.
\end{resumen}

\section{Cómo se encuentran los dos puntos}\label{sec:deteccion}

Cada cuadro del tramo, recortado a la región, pasa por seis pasos. Se procesan bloques de muchos cuadros a la vez, con
operaciones de NumPy sobre arreglos de forma (cuadros, alto, ancho), sin recorrer los píxeles uno por uno.

\begin{enumerate}
\item \textbf{Fondo y ruido.} El fondo es la mediana del recorte y el ruido, $1{,}4826$ veces la desviación absoluta
mediana (MAD, sección~\ref{sec:robusta}). Si hay maestros de calibración, el cuadro se corrige antes como
$(\text{cuadro} - \text{dark})/\text{flat}$.
\item \textbf{Píxeles con luz.} Se suaviza la imagen con un promedio de 3~×~3 píxeles, para que un píxel caliente aislado
no parezca una estrella, y se toman los píxeles que superan
$\max(0{,}15\times\text{pico},\ \tfrac43\times\text{ruido})$.
\item \textbf{Dos grupos.} Esos píxeles se reparten en dos grupos con \emph{k-medias} ($k = 2$) ponderado por el
brillo: el primer grupo arranca en el píxel más brillante y el segundo en el píxel brillante más lejano de él; cada píxel
se asigna al centro más cercano y los centros se recalculan, hasta ocho veces o hasta que no se muevan. Así se separan bien
puntos pequeños, grandes, en forma de anillo (desenfocados) o partidos en pedazos por la turbulencia.
\item \textbf{Centroides.} El centro de cada punto es el promedio de las posiciones de sus píxeles, ponderado por su
brillo. Los píxeles a más de $3r + 3$ píxeles del centro (con $r$ el radio de giro del grupo) se descartan y el centro se
recalcula: así un reflejo o una estrella vecina no lo arrastran. La precisión es de centésimas de píxel.
\item \textbf{Validación.} El cuadro cuenta solo si los dos grupos son dos puntos de verdad: los dos con un pico de más de
seis veces el ruido, el más débil con al menos el 12\,\% de la luz del otro, separados más de 2,2 veces su tamaño, con una
\textbf{caída de luz} entre ellos (en el punto medio, menos del 60\,\% del pico menor) y sin tocar el borde del recorte.
Si se ve un solo punto, no se parte en dos.
\item \textbf{Orientación.} Los dos puntos pueden intercambiar su orden de un cuadro a otro (cambia cuál es el más
brillante). TurbuLENS orienta todos los vectores según el eje principal de la nube de separaciones, de modo que la
componente longitudinal tenga siempre el mismo sentido.
\end{enumerate}

\begin{figure}[H]
\centering
\begin{tikzpicture}[font=\sffamily\scriptsize]
  \foreach \i/\t in {0/Cuadro,1/Píxeles con luz,2/Dos grupos,3/Centroides}{
    \begin{scope}[shift={(\i*3.6,0)}]
      \fill[tinta] (0,0) rectangle (3,3);
      \node[anchor=north,font=\sffamily\small] at (1.5,-0.1) {\t};
    \end{scope}}
  % cuadro con ruido
  \foreach \k in {1,...,70}{\pgfmathsetmacro{\xx}{0.1+2.8*rnd}\pgfmathsetmacro{\yy}{0.1+2.8*rnd}
    \fill[white!30!gris] (\xx,\yy) circle (0.02);}
  \shade[inner color=white,outer color=tinta] (1.0,1.9) circle (0.42);
  \shade[inner color=white!80!naranja,outer color=tinta] (2.0,1.05) circle (0.38);
  % umbral
  \begin{scope}[shift={(3.6,0)}]
    \foreach \dx/\dy in {0/0,0.12/0,-0.12/0,0/0.12,0/-0.12,0.12/0.12,-0.12/-0.12,0.12/-0.12,-0.12/0.12,0.24/0,-0.24/0,0/0.24,0/-0.24}{
      \fill[white] (1.0+\dx-0.05,1.9+\dy-0.05) rectangle ++(0.1,0.1);
      \fill[white] (2.0+\dx-0.05,1.05+\dy-0.05) rectangle ++(0.1,0.1);}
  \end{scope}
  % grupos
  \begin{scope}[shift={(7.2,0)}]
    \foreach \dx/\dy in {0/0,0.12/0,-0.12/0,0/0.12,0/-0.12,0.12/0.12,-0.12/-0.12,0.12/-0.12,-0.12/0.12,0.24/0,-0.24/0,0/0.24,0/-0.24}{
      \fill[acento!60!white] (1.0+\dx-0.05,1.9+\dy-0.05) rectangle ++(0.1,0.1);
      \fill[naranja] (2.0+\dx-0.05,1.05+\dy-0.05) rectangle ++(0.1,0.1);}
    \draw[white,dashed] (0.3,0.4) -- (2.7,2.6);
  \end{scope}
  % centroides
  \begin{scope}[shift={(10.8,0)}]
    \shade[inner color=white,outer color=tinta] (1.0,1.9) circle (0.42);
    \shade[inner color=white!80!naranja,outer color=tinta] (2.0,1.05) circle (0.38);
    \draw[acento!50!white,line width=0.8pt] (1.0,1.9) -- (2.0,1.05);
    \draw[white,line width=0.8pt] (0.88,1.9) -- (1.12,1.9) (1.0,1.78) -- (1.0,2.02);
    \draw[white,line width=0.8pt] (1.88,1.05) -- (2.12,1.05) (2.0,0.93) -- (2.0,1.17);
    \node[text=white] at (1.75,1.7) {$\vec d$};
  \end{scope}
\end{tikzpicture}
\caption{De un cuadro al vector $\vec d$ entre los dos centroides.}
\end{figure}

Después de procesar todos los cuadros, los que tienen una separación a más de 5 desviaciones robustas de la mediana se
descartan como atípicos (sección~\ref{sec:robusta}). Con los que quedan se calculan $\sigma_\ell$ y $\sigma_t$.

\begin{nota}[Probado con valores conocidos.]
La detección se prueba con videos sintéticos en los que sabemos dónde está cada punto: con puntos pequeños, grandes, en
anillo y con ruido, el error típico del centroide es de 0,02 a 0,05 píxeles. En un video real de puntos grandes y
desenfocados, el programa con el que empezó este proyecto, que usaba un umbral fijo, encontraba un solo objeto en 163 de
462 cuadros; con los dos grupos, 457 de los 462 cuadros son válidos.
\end{nota}

\section{Media, desviación estándar y error}

Para $N$ datos $x_i$ (las separaciones, o sus componentes):
\begin{align}
\bar x &= \frac{1}{N}\sum_{i=1}^{N} x_i, &
s &= \sqrt{\frac{1}{N-1}\sum_{i=1}^{N}(x_i - \bar x)^2}, &
\text{EEM} &= \frac{s}{\sqrt N}.
\end{align}
La desviación estándar usa $N-1$ (la versión muestral) y el error estándar de la media (EEM) dice cuánto variaría
$\bar x$ si se repitiera la medición. Para el DIMM lo que importa es $s$, no $\bar x$: la separación media depende de
cuánto se desenfocó y no dice nada de la atmósfera; su variación sí.

\textbf{Histogramas.} El número de intervalos sigue la regla de Freedman–Diaconis: ancho
$h = 2\,\text{RIC}/N^{1/3}$, con RIC el rango intercuartílico, limitado a entre 8 y 60 intervalos.

\section{Estadística robusta}\label{sec:robusta}

La media y la desviación estándar tienen un defecto: \textbf{un solo dato muy lejano las arrastra}. En la lección 7 del
modo educativo se cuenta un caso real: 68 cuadros mal medidos, con distancias cercanas a cero, entre 8\,484 buenos, casi
duplicaron la desviación estándar, y con ella el seeing.

La \textbf{mediana} y la \textbf{desviación absoluta mediana} no tienen ese problema:
\begin{equation}
\text{MAD} = \operatorname{mediana}\left(|x_i - \operatorname{mediana}(x)|\right),
\qquad \sigma_\text{robusta} = 1{,}4826 \times \text{MAD}.
\end{equation}
El factor 1,4826 hace que, para datos normales, $\sigma_\text{robusta}$ coincida con la desviación estándar. Hasta la
mitad de los datos pueden ser basura sin que la mediana y la MAD se muevan mucho.

TurbuLENS descarta los cuadros con
\begin{equation}
|x_i - \operatorname{mediana}(x)| > 5\,\sigma_\text{robusta} ,
\end{equation}
y repite el recorte hasta tres veces o hasta que no salga ningún dato más. El umbral de 5 es deliberadamente generoso:
en una distribución normal, solo uno de cada 1,7 millones de datos buenos estaría tan lejos, así que no se recortan las
colas reales de la turbulencia, solo los errores.

\section{Correlación}

El \textbf{coeficiente de correlación de Pearson} entre dos variables medidas en los mismos cuadros es
\begin{equation}
r = \frac{\sum_i (x_i - \bar x)(y_i - \bar y)}{\sqrt{\sum_i (x_i - \bar x)^2\,\sum_i (y_i - \bar y)^2}},\qquad -1\le r\le 1 .
\end{equation}
Vale 1 si los puntos $(x_i, y_i)$ están sobre una recta que sube, $-1$ si están sobre una que baja, y 0 si no hay
relación lineal. $r^2$ es la fracción de la variación de una variable que «explica» la otra.

\section{Los tres controles de calidad}\label{sec:controles}

Con los mismos datos, TurbuLENS calcula tres números que dicen si lo medido se comporta como turbulencia. Aparecen debajo
del diagrama de movimiento diferencial de \menu{Resultados} y se explican, con los valores de cada medición, en la lección 8
del modo educativo.

\begin{description}[style=nextline,leftmargin=1.2em,font=\bfseries\color{acentoosc}]
\item[$r$ longitudinal–transversal]
La turbulencia mueve las dos componentes de forma independiente, así que lo esperado es un valor cercano a 0. TurbuLENS
considera bueno $|r| < 0{,}2$. Un $|r|$ grande, positivo o negativo, delata algo que no es turbulencia: vibración o viento
en una sola dirección, cambios de foco, una línea de base mal estimada o puntos confundidos. El signo solo depende de cómo
quedaron los ejes.
\item[$r$ entre los dos puntos]
Mide cuánto se mueven juntas las dos imágenes, por separado en cada componente, después de quitar la deriva lenta (un
promedio móvil de un segundo, que saca el desplazamiento del seguimiento). Cerca de 1: domina el movimiento común, como
la vibración del telescopio; el DIMM lo cancela al medir la diferencia, así que el seeing sigue valiendo, pero si es muy
grande puede emborronar los puntos. Entre 0 y 1: lo normal, porque los dos orificios ven parte de la misma turbulencia.
Cerca de 0: cada orificio ve turbulencia independiente. Negativo: los puntos se acercan y se alejan de forma opuesta
(cambios de foco o errores de detección).
\item[$\sigma_t/\sigma_\ell$ medido frente al teórico]
El teórico es $\sqrt{K_t/K_\ell}$, que depende solo de $B/D$. Si el medido se aparta más de un 25\,\% del teórico,
TurbuLENS lo señala.
\end{description}

\begin{center}
\begin{tikzpicture}[font=\sffamily\scriptsize]
  \foreach \i/\t/\c in {0/{Turbulencia: nube redonda, $r\approx0$}/acento,1/{Vibración en diagonal: $|r|$ grande}/naranja}{
    \begin{scope}[shift={(\i*6.4,0)}]
      \draw[gris!50] (-2.2,0) -- (2.2,0) (0,-1.6) -- (0,1.6);
      \node[anchor=north,font=\sffamily\small] at (0,-1.75) {\t};
    \end{scope}}
  \pgfmathsetseed{7}
  \foreach \k in {1,...,160}{
    \pgfmathsetmacro{\a}{0.55*(rnd+rnd+rnd+rnd-2)*1.6}\pgfmathsetmacro{\b}{0.5*(rnd+rnd+rnd+rnd-2)*1.6}
    \fill[acento,opacity=0.55] (\a,\b) circle (0.035);
    \pgfmathsetmacro{\u}{0.55*(rnd+rnd+rnd+rnd-2)*1.6}\pgfmathsetmacro{\v}{0.18*(rnd+rnd+rnd+rnd-2)*1.6}
    \fill[naranja,opacity=0.6] ({6.4+0.8*\u-0.6*\v},{0.6*\u+0.8*\v}) circle (0.035);}
  \node[text=gris] at (2.0,1.45) {$t$}; \node[text=gris] at (2.25,-0.18) {$\ell$};
\end{tikzpicture}
\end{center}

\section{La desviación estándar en el tiempo}

La gráfica \emph{Dispersión segundo a segundo} de \menu{Resultados} divide el tramo en ventanas de un segundo y calcula
la desviación estándar de la separación en cada una. Si la turbulencia fuera estable, la línea sería casi plana; en la
práctica sube y baja un poco porque el seeing cambia de verdad. Lo que hay que mirar son los \textbf{saltos aislados}:
un golpe al telescopio o una ráfaga de viento sobre el tubo dura uno o dos segundos y dispara $\sigma$ en esas ventanas.
Dejarlos fuera del tramo da un seeing más honesto.
